\originalchapter{8.04 作业3：对易、相位与概率流}
以下按MIT原件中文转述, 解答均为\textbf{AI独立解答}; 保留原件的改编信用, 依完整MIT任务独立推导. 来源为Spring 2016 Assignment 3, 2016年2月18日发布, 2月25日17时截止.
\originalsection{第1题：对易子的代数}
\begin{exercise}对线性算符$A,B,C$证明$[A,BC]=[A,B]C+B[A,C]$及$[AB,C]=A[B,C]+[A,C]B$, 再证Jacobi恒等式. 求$[x^n,p]$、$[x,p^n]$（$n>0$整数）及$[xp,x^2]$、$[xp,p^2]$.\end{exercise}
\begin{solution}第一式右侧展开为$ABC-BAC+BAC-BCA=ABC-BCA$, 第二式同法相邻消项. 将$[A,[B,C]]$等三个嵌套式各展开成四项, 每个有序乘积各出现一次正号与一次负号, 总和为0. 从$[x,p]=\ii\hbar I$出发反复用乘积式, 给$[x^n,p]=\ii\hbar nx^{n-1}$及$[x,p^n]=\ii\hbar np^{n-1}$. 最后$[xp,x^2]=x[p,x^2]=-2\ii\hbar x^2$, $[xp,p^2]=[x,p^2]p=2\ii\hbar p^2$. 无界算符在共同光滑定义域计算; 形式等式不自动解决定义域问题.\end{solution}
\originalsection{第2题：驻相与Gaussian中心}
\begin{exercise}$\Psi(x)=\int_{-\infty}^{\infty}\Phi(k)\e^{\ii kx}\dd k$, 分别取$\Phi=\e^{-L^2(k-k_0)^2}$及再乘$\e^{-\ii kx_0}$, 用相位预测峰位置并精确积分核对, $L>0$.\end{exercise}
\begin{solution}窄幅包络内相位变化最慢的位置分别满足$x=0$、$x-x_0=0$. 令$q=k-k_0$完成平方, 第一种$\Psi=\sqrt\pi\,\e^{\ii k_0x}\e^{-x^2/(4L^2)}/L$, 第二种把$x$全部换成$x-x_0$. 模分别为$\sqrt\pi\e^{-x^2/(4L^2)}/L$及其平移, 精确峰与预测一致. 这里“驻相”指把相位对积分变量的导数置零; 线性相位只在这些特定$x$处全区间同相, 不是寻找包络对$k$的极值.\end{solution}
\originalsection{第3题：Galilei变换}
\begin{exercise}自由方程$\ii\hbar\partial_t\Psi=-\hbar^2\partial_x^2\Psi/(2m)$在$x'=x-vt,t'=t$下应仍成立. 求普适因子$f(x,t)$使$\Psi'(x',t')=f(x,t)\Psi(x,t)$, 并核平面波的动量和能量.\end{exercise}
\begin{solution}链式法则在固定$x'$时给$\partial_{t'}=\partial_t+v\partial_x$. 设$f=\e^{\ii(Ax+Bt)}$, 代入并比较任意$\Psi$的$\partial_x\Psi$系数得$A=-mv/\hbar$, 比较$\Psi$系数得$B=mv^2/(2\hbar)$. 所以$f=\exp[-\ii mvx/\hbar+\ii mv^2t/(2\hbar)]$, 可另乘整体常相位. 原平面波$\e^{\ii(px-Et)/\hbar}$变成$\e^{\ii[(p-mv)x'-(E-pv+mv^2/2)t']/\hbar}$. 于是$p'=p-mv$, $E'=p'^2/(2m)$, 符合新参考系的自由能量; 只换坐标却漏$f$将得到错误动量.\end{solution}
\originalsection{第4题：三维连续性}
\begin{exercise}从三维实势Schrödinger方程及$\rho=|\Psi|^2$推$\partial_t\rho+\boldsymbol\nabla\cdot\boldsymbol J=0$, 求$\boldsymbol J$.\end{exercise}
\begin{solution}把方程乘$\Psi^*$, 共轭方程乘$\Psi$后相减, 实势两项抵消, 得$\partial_t\rho=\ii\hbar(\Psi^*\nabla^2\Psi-\Psi\nabla^2\Psi^*)/(2m)$. 用散度乘积式把括号写为$\nabla\cdot(\Psi^*\nabla\Psi-\Psi\nabla\Psi^*)$, 所以$\boldsymbol J=\hbar\operatorname{Im}(\Psi^*\nabla\Psi)/m$. 在体积上积分给$\dot P=-\oint\boldsymbol J\cdot\dd\boldsymbol a$; 三维$\rho$量纲$L^{-3}$, $J$量纲$L^{-2}T^{-1}$.\end{solution}
\originalsection{第5题：两个展宽波包的内积}
\begin{exercise}
本题保留原件Merzbacher署名. 初始两个各自可归一窄波包相距很远, $\Psi(0)=\Psi_1(0)+\Psi_2(0)$, $\gamma(t)=\int\Psi_1^*(x,t)\Psi_2(x,t)\dd x$, $|\gamma(0)|$很小. 波包展宽并在位置上重叠后, $|\gamma|$会否增加?
\end{exercise}
\begin{solution}
两波包按同一个自伴Hamilton算符演化. 直接求导给$\dot\gamma=(\ii/\hbar)\langle H\Psi_1,\Psi_2\rangle-(\ii/\hbar)\langle\Psi_1,H\Psi_2\rangle=0$, 所用定义域和边界条件使两项相等. 等价地$\langle U\Psi_1,U\Psi_2\rangle=\langle\Psi_1,U^\dagger U\Psi_2\rangle=\gamma(0)$. 因此复内积本身保持. 位置概率密度可能越来越重叠, 但积分中的相对相位造成抵消, 几何重叠不是态内积的增大.
\end{solution}
\originalsection{第6题：三种波函数的流}
\begin{exercise}计算$t=0$时(a)$A\e^{\gamma x}$（$\gamma$实）, (b)$N(x)\e^{\ii S(x)/\hbar}$（$N,S$实）, (c)$A\e^{\ii kx}+B\e^{-\ii kx}$的概率流.\end{exercise}
\begin{solution}(a)$\Psi^*\Psi'=\gamma|A|^2\e^{2\gamma x}$为实, 故$J=0$; 不归一的表达只作局部流计算. (b)$\Psi^*\Psi'=NN'+\ii N^2S'/\hbar$, 故$J=N^2S'/m$. (c)交叉项$\ii k[B^*A\e^{2\ii kx}-A^*B\e^{-2\ii kx}]$为实, 只剩$J=\hbar k(|A|^2-|B|^2)/m$. 密度含干涉, 流却没有左右交叉贡献.\end{solution}
